% Framework non-equivalence table (reviewer R2).
\begin{table*}[t]
\centering
\caption{The three instruments are \emph{not} equivalent assessment objects. A shared
technical artifact may be relevant to all three, but their legal status, unit of assessment,
responsible actor, and assurance mechanism differ. In particular, ISO/IEC~42001 certification
certifies an organization's management system, not that an individual AI system is EU~AI~Act
compliant, and the NIST AI RMF is voluntary guidance, not law. \textbf{Takeaway:} the three
regimes overlap heavily in \emph{evidence} but differ in legal force, unit of assessment,
responsible actor, and assurance mechanism; shared evidence must not be read as legal
equivalence.}
\label{tab:frameworks}
\footnotesize
\setlength{\tabcolsep}{4pt}
\begin{tabular}{p{0.16\textwidth}p{0.25\textwidth}p{0.27\textwidth}p{0.24\textwidth}}
\toprule
Dimension & NIST AI RMF 1.0 & EU AI Act (Reg.\ (EU) 2024/1689) & ISO/IEC 42001:2023\\
\midrule
Legal status & Voluntary guidance & Binding law (phased application) & Voluntary, certifiable standard\\
Object assessed & AI risk-management \emph{practice} of an organization & The \emph{high-risk AI system}, per its intended purpose & The organization's \emph{AI management system} (AIMS)\\
Responsible actor & Any AI actor (adopter) & Provider (with deployer duties) & The organization seeking certification\\
Assurance mechanism & Self-directed; no external conformity assessment & Conformity assessment, self- or notified-body (Art.~43) & Third-party certification audit\\
Lifecycle role & GOVERN/MAP/MEASURE/MANAGE, lifecycle-wide & Pre-market \emph{and} post-market obligations (Arts.~9--15, 72) & Plan--Do--Check--Act management cycle\\
Role of technical evidence & Populates the MEASURE function & One input among QMS, documentation, and the CA route & Evidence that AIMS controls operate as intended\\
\bottomrule
\end{tabular}
\end{table*}
